\documentclass[11pt]{article}
\usepackage[margin=1in]{geometry}
\usepackage{amsmath}
\usepackage{booktabs}
\usepackage{xcolor}
\usepackage{hyperref}
\usepackage{parskip}
\usepackage{listings}
\lstset{basicstyle=\ttfamily\small, breaklines=true, columns=fullflexible}

\definecolor{warnbg}{RGB}{255,243,205}
\newcommand{\warnbox}[1]{%
  \par\noindent\colorbox{warnbg}{\parbox{\dimexpr\linewidth-2\fboxsep}{\textbf{Caveat:} #1}}\par\vspace{4pt}
}

\title{EpiGen: where the wall-clock time actually goes}
\author{}
\date{\today}

\begin{document}
\maketitle

\section*{Why this document exists}

A run was timed end to end on the whole-loop-on-Modal path
(\texttt{use\_modal\_mcmc=True}) to answer ``what's taking so long.'' This is
that measurement, what it does and doesn't tell us, and the concrete levers
available to cut it down.

\tableofcontents

\section{The measured run}

\begin{center}
\begin{tabular}{@{}ll@{}}
\toprule
Config & Value \\
\midrule
Scaffold & 193L (lysozyme, 129 aa) \\
Edit & position 60, \texttt{WHSPRAL} (7 residues) \\
Window & full protein minus the edit (122 positions) \\
\texttt{num\_starting\_points} $\times$ \texttt{chains\_per\_start} & $20 \times 2 = 40$ chains \\
\texttt{steps} & 20 (upper bound; chains freeze early) \\
\texttt{candidate\_num} & 5 \\
\texttt{use\_modal\_mcmc} & \texttt{True} (whole search loop runs inside one Modal call) \\
\bottomrule
\end{tabular}
\end{center}

\textbf{Total wall time: 1075.7s (\textasciitilde 18.0 min).} Result:
5 candidates, \texttt{wt\_score=720.210}, \texttt{wt\_nt\_source=codonfm}
(193L has no verified GenBank CDS for this sequence -- Evo2's nt sequence
came from the CodonFM fallback, not a real deposited coding sequence; see
\texttt{oracle.nt\_resolution}).

\section{The two-bucket breakdown}

The client-visible log only prints a timer around each \emph{separate} Modal
dispatch (``Connecting to modal container: 0it [XXs, ?it/s]''). Summing the
non-zero ones gives the overhead \emph{outside} the search loop itself; the
remainder is time spent inside the one, single Modal call that runs the
entire MCMC search server-side.

\begin{center}
\begin{tabular}{@{}lrr@{}}
\toprule
Bucket & Time & Share \\
\midrule
Visible container-connect overhead (sum of non-zero ``Connecting...'' timers) & \textasciitilde 204s & 19\% \\
Inside the single \texttt{run\_mcmc\_search\_remote} call (no separate connect events) & \textasciitilde 872s & 81\% \\
\midrule
Total & 1075.7s & 100\% \\
\bottomrule
\end{tabular}
\end{center}

The 204s bucket is the one-off calls \emph{around} the search: the initial
fold, ProteinMPNN warm-start generation (\texttt{propose\_compensatory\_mutations}),
the pre/post-search ESM2+ProteinMPNN oracle-sanity-check batch calls, and
refold+TM-align on the top candidates. This is exactly the overhead
\texttt{use\_modal\_mcmc=True} was built to shrink for the \emph{search
itself} (collapsing per-round laptop$\leftrightarrow$Modal round trips into
one container-to-container loop) -- and it did: none of the 20 MCMC rounds
shows up here as a separate connect event, because they're not separate
dispatches anymore.

The 872s bucket -- the large majority of the wall time -- is therefore
\emph{not} network/dispatch overhead. It's real compute time inside one
container: up to 20 rounds, each round scoring up to 40 chains' proposals
with three experts (ESM2, ProteinMPNN, Evo2).

\warnbox{This split is derived from the client-side log's dispatch timers,
not from any instrumentation inside \texttt{run\_mcmc\_search}/
\texttt{run\_mcmc\_search\_remote} itself. It tells us \emph{overhead vs.\
everything else}, not which of the three experts (or which round) dominates
the 872s. See \S\ref{sec:next} for what real per-expert numbers would take.}

\section{Why the inner 872s is plausibly real compute, not something broken}

Three things compound inside that one call:

\begin{itemize}
  \item \textbf{Round count is worst-case, not average-case.} The round loop
    only exits early once \emph{every} chain has frozen (its own score
    exceeded \texttt{wt\_score}). A handful of slow-to-converge chains out of
    40 can force close to the full 20 rounds even if most chains froze
    much earlier.
  \item \textbf{Every active round pays three full batched model calls} --
    \texttt{position\_scores\_esm2\_batch}, \texttt{position\_scores\_proteinmpnn\_batch},
    and \texttt{window\_log\_prob\_batch} (Evo2) -- each over every still-active
    chain's proposal.
  \item \textbf{Evo2's call is the heaviest of the three per-token.} As of
    today's window-restriction fix, it runs with \texttt{return\_logits=True} --
    a full \texttt{(seq\_len, vocab\_size)} logits tensor per sequence, not
    just a scalar \texttt{avg\_log\_likelihood} -- so it's doing meaningfully
    more work per call than the old whole-sequence-average version, in
    exchange for being scoped correctly (window-only, matching ESM2/ProteinMPNN)
    instead of diluted over the whole sequence.
\end{itemize}

None of this is a bug; it's the cost of the fix instead of a symptom of
it -- the network overhead is gone, what's left is the model compute the
search was always going to need for a 40-chain, up-to-20-round, 3-expert
search.

\section{Levers to actually cut wall time}

\begin{center}
\begin{tabular}{@{}lp{0.65\textwidth}@{}}
\toprule
Lever & Effect \\
\midrule
Lower \texttt{steps} & Directly caps worst-case rounds; the freeze mechanism
  already exits early for chains that converge, so this mostly trims the
  slow tail. \\
Fewer chains (\texttt{num\_starting\_points} $\times$ \texttt{chains\_per\_start}) &
  Shrinks every round's batch size linearly; 40 chains is on the high end of
  what's been run this session. \\
Raise \texttt{weight\_evo2} or tune \texttt{temperature} & Changes how fast
  chains cross the freeze threshold -- may shorten the effective round count
  without touching \texttt{steps} itself, but changes what's being optimized. \\
Drop Evo2 (\texttt{use\_evo2=False}) & Removes the heaviest per-round call
  entirely -- back to a 2-expert search, faster but loses the DNA-level
  signal. \\
\bottomrule
\end{tabular}
\end{center}

\section{What would give a real per-expert breakdown}
\label{sec:next}

Right now there is no timestamp recorded around the individual ESM2 /
ProteinMPNN / Evo2 calls \emph{inside} \texttt{run\_mcmc\_search}, on either
the laptop-orchestrated or whole-loop-on-Modal path. Getting genuine
per-expert numbers (not just ``inside vs.\ outside the big call'') would mean
adding \texttt{time.time()} timers around each of the three batched calls in
\texttt{oracle/mcmc.py}'s round loop, aggregating them (e.g.\ into the same
checkpoint progress payload \texttt{oracle/checkpoint.py} already writes),
and re-running. Not done here -- this document is the coarse split available
from the existing log output, not a profiling pass.

\end{document}
